\section{Metodología experimental}\label{sec:metodologia}

\subsection{Preguntas e hipótesis}
\begin{table}[H]
\centering\small
\caption{Preguntas de evaluación e hipótesis de la Entrega 1.}
\begin{tabularx}{\textwidth}{cL{6.0cm}Y}
\toprule
\# & Pregunta & Hipótesis \\
\midrule
Q1 & ¿Qué estrategia atiende más vueltas y obtiene menor $\Lmax$? & H1: sin tráfico LS y LPT coinciden (las vueltas de una ruta duran lo mismo); con tráfico LS supera a LPT. \\
Q2 & ¿Cómo afecta el tráfico a la cantidad de vueltas atendidas? & H2: T1 tiene un efecto moderado; T2 reduce marcadamente la atención. \\
Q3 & ¿Qué rutas son más difíciles de programar? & H3: las vueltas no atendidas se concentran en las rutas de mayor saturación $\rho$. \\
\bottomrule
\end{tabularx}
\end{table}

\subsection{Experimento E1: instancias, variables y métricas}
\begin{itemize}
  \item \textbf{Instancias:} las 35 rutas experimentales (4\,733 vueltas y 1\,080 buses).
  \item \textbf{Variables:} algoritmo \{LS, LPT\} $\times$ perfil de tráfico \{T0, T1, T2\}, es decir, 6 configuraciones por ruta.
  \item \textbf{Métricas principales} (función objetivo): vueltas no atendidas $U$ y carga máxima $\Lmax$.
  \item \textbf{Métricas complementarias:} brecha de $\Lmax$ respecto de la cota inferior de las vueltas atendidas, desbalance $\Lmax - L_{\min}$, sobrecosto por tráfico $\bigl(\sum_j p_j(s_j) - \sum_j \pbar_j\bigr) / \sum_j \pbar_j$, tiempo de ejecución y violaciones detectadas por el validador.
\end{itemize}
La metodología global del proyecto prevé, en etapas posteriores, la comparación con los mecanismos de referencia en instancias pequeñas, el análisis de escalabilidad y la evaluación de la robustez ante la variabilidad del tráfico. Esos análisis no forman parte de esta entrega.

\subsection{Entorno y procedimiento reproducible}
\begin{itemize}
  \item Equipo: computadora del grupo con Windows (PowerShell) y un entorno virtual \texttt{.venv}. Software: Python $\ge$ 3.10 con las dependencias de \texttt{requirements.txt} (pandas, openpyxl, matplotlib y pytest); el paquete \texttt{bus\_sched} se instala en modo editable (\texttt{pip install -e .}). Los resultados y capturas de este informe provienen de ese equipo.
  \item Procedimiento (desde la raíz del repositorio):
\begin{lstlisting}
python -m bus_sched build              # genera las 35 instancias JSON
python experiments/run_experiments.py  # E1 -> results/E1_35_rutas.csv, E1_resumen.csv y figuras
python experiments/build_excel.py      # Excel de la entrega
python experiments/export_latex.py     # tablas y figuras de este informe
\end{lstlisting}
  \item Los algoritmos son deterministas: repetir el procedimiento produce los mismos resultados (salvo los tiempos de ejecución).
\end{itemize}
